\scanpage{21}
% L1–2
\uline{et par suite aussi sur le faisceau} $\boxed{\underline S(\overline X/X_0)}$ \uline{des germes de sections de} $\overline X$ \uline{sur} $X_0$

% L3–8
On peut aussi considérer le faisceau $\underline{\operatorname{Is}}_T^{X_0}(\deleted{X}\xi,\deleted{\xi}X)$, qui est un faisceau \uline{principal homogène} (\unsure{référence|v. finesse|références|vérifiant} \ldots) sous $\underline{\operatorname{Aut}}_T^{X_0}(X)$ (opérant à gauche) et sous $\underline{\operatorname{Aut}}_T^{X_0}(\xi)$ (\rule{2.5cm}{0.35pt} droite). D'ailleurs on a une section privilégiée \inserted{$s_0$} de $\underline S(\overline{\xi}/\xi_0)$

% L9–11
et on en déduit un hom de faisceaux
\[
u\longmapsto u(s_0)
\]
\[
\overline{\underline{\operatorname{Is}}}_T^{X_0}(\xi,X)\longrightarrow\underline S(\overline X/X_0)
\]

% L12–15, vertical rule to the left
\uline{Proposition} L'homomorphisme précédent est un isomorphisme de faisceaux, \deleted{\unsure{de groupes à}} \inserted{compatible avec le} faisceau de groupes d'opérateurs $\underline{\operatorname{Aut}}_T^{X_0}(X)$ (à gauche).

% L16–17
\uline{Corollaire 1} $\underline S(\overline X/X_0)$ est \uline{principal homogène} sous le faisceau $\underline{\operatorname{Aut}}_T^{X_0}(X)$

% L17–19, left
% slanted
\margin{Ceci vaut $=$ si $\underline{\Lambda}$ \unsure{---} $\longleftrightarrow$ \unsure{qu'une|par une} $\underline{O}_S$-Algèbre, \unsure{mais} \unsure{---} sans $\underline{O}_S$-augmentation, cf. Greenberg pour la définition de $\overline X$}
% L18–20
\uline{Corollaire 2} $\underline{\operatorname{Aut}}_T^{X_0}(\xi)$ opère de façon naturelle \inserted{à droite} sur $\underline S(\overline X/X_0)$, qui devient un \uline{faisceau principal hom.} à droite sous ce faisceau de groupes \deleted{\unsure{à droite}}
% archive sheet number: 20

